% !PRINTED 23
\weberpage{23}
If one sets
\weq{9}{\frac{J}{\Delta^2} = 3z, \quad dz = \frac{1}{3}\frac{\Delta dJ-2Jd\Delta}{\Delta^3},}
then from (8) there results
\[
K^2 = 27\Delta^6(4z^3+12Sz-T)
\]
and from (7)
% [sic] The source prints 1/\sqrt{3}; equations (7)--(9) algebraically give \sqrt{3} (up to the branch sign).
\weq{10}{du = \frac{1}{\sqrt3}\frac{dz}{\sqrt{4z^3+12Sz-T}},}
and this is the \gesperrt{Weierstrass} normal form for $g_2=-12S$,
$g_3=T$.

\begin{center}
\S\,7. \gesperrt{Elliptic Space Curves of the Fourth Order.}
\end{center}

Algebraic functions of one variable can also be represented by
equations between several variables, provided one increases the
number of equations correspondingly. If one takes, for example,
three variables $x$, $y$, $z$ and lets two equations
\weq{1}{\varphi(x,y,z) = 0, \qquad \psi(x,y,z) = 0}
hold between them, then out of these two equations one can, for
example, eliminate $x$ and obtains an equation between $y$ and $z$ by
which $y$ is defined as an algebraic function of $z$. If certain
exceptional cases are excluded, $x$ and every rational function of
$x$, $y$, $z$ can then be represented rationally by $y$ and $z$. The
integrals of the form
\weq{2}{\int F(x,y,z)\,dz,}
in which $F$ denotes a rational function, then belong to the
algebraic figure represented by (1). If one takes
$x$, $y$, $z$ as Cartesian coordinates in space, then (1) represents
a space curve as the intersection of two surfaces $\varphi=0$,
$\psi=0$, and the integrals (2) belong to this space curve.
The curve is again called \gesperrt{elliptic} if these integrals are elliptic.

We wish to apply these considerations to the \gesperrt{space curves of the
fourth order of the first species}, that is, to the curves of the
fourth order which can be represented as the complete intersection of
two surfaces of the second degree.
